\section{Conclusion}
\label{sec:conclusion}

We built a general, microscopic, governance-generative agent-based model
from elementary mechanisms addressing cooperation, ecological viability,
information, enforcement, delegation, accountability, mobility, and
interpatch coordination, and, across four linked milestones, searched its
outcome space, blind to any expected answer, for recurrent institutional
phenotypes and for the causal role of history in producing them. The
synchronic substrate (Level 1.2/1.2b) is generatively sufficient ---
producing institutions and heterogeneous institutional outcomes from
purely local rules --- but not, by itself, sufficient to compress those
outcomes into a small number of discrete, well-separated forms; both a
24-run and a 576-run/384-run independently replicated search returned a
\texttt{continuum} verdict. Extending that substrate with historically
structured pressure, sequencing, entrenchment, and legitimacy (Level
1.3/1.3b) did not change this qualitative picture, but it did add
something the synchronic substrate could not: direct, sham-validated
causal evidence that history matters. Erasing, reversing, or
exposure-equalizing a simulated jurisdiction's own realized pressure
history produces a real, persistent, non-decaying divergence from that
same jurisdiction with its history intact, validated against a control
proven to introduce zero spurious divergence on its own --- the strongest
causal evidence for path dependence this project has produced. At the same
time, a purpose-designed, preregistered discreteness test, reusing the
synchronic substrate's own frozen classification rule unmodified, rejected
discrete institutional families more decisively for the historically
extended substrate than for the synchronic one: no tested cluster count
cleared even an entry-level silhouette floor, cross-algorithm agreement was
weak, and the leading partition did not reproduce across ecological
strata. History, in this model, is a real, causal, persistent conditioning
operator on where a governance configuration sits and how it got there ---
without collapsing the space of configurations into a fixed small number
of discrete attractors.

We have argued throughout that these two findings are not in tension. The
dissertation's own $S$/$Z$ coordinate space remains the central,
interpretively load-bearing frame for this paper's own results: the
model's independently constructed $S_{\text{obs}}$/$Z_{\text{obs}}$
composites show a real, if modest and weighting-sensitive, positive
coupling, weakly corroborating the dissertation's own finding that state
capacity and institutional constraint are related but distinct dimensions
of governance. What the historically extended substrate adds is not a
third discrete axis but a richer account of movement through that same
$S$/$Z$ space --- accumulated pressure history, entrenchment, and
legitimacy that jointly explain why a configuration sits where it does and
how it is likely to move next, formalized here as
$\mathcal{G}_t=[S_t,Z_t,H_t,E_t,L_t,P_t,R_t,C_t,\ldots]$ and
$G_t=f(S_t,Z_t;H_{0:t},E_t,L_t,P_t,\ldots)$. Three complementary lines of
evidence triangulate toward this account without duplicating one another:
the dissertation's own SEM identification of a two-dimensional coordinate
space from real survey data; this project's demonstration that a
microscopic mechanism set, encoding none of the dissertation's own
variables or labels, is generatively sufficient to populate an analogous
space and to move heterogeneously within it under historical pressure; and
a sham-validated causal intervention design, structurally unavailable to
cross-sectional survey analysis, showing that history's effect on
simulated institutional trajectory is causal rather than merely
associational. None of these three lines is sufficient on its own, and we
have been correspondingly careful, throughout this paper's own
terminological discipline, not to let a statistical component count stand
in for a species count, a recurrent phenotype for a dynamical attractor,
or a positive generalization gap against a weak null for genuine
discreteness.

The dissertation's three empirically observed governance trajectories
remain a legitimate, unaffected empirical finding at their own analytical
level, and this paper's own results neither confirm nor refute them
directly. What this paper's evidence does support is a specific,
bounded reading: those three trajectories are best understood as
historically and institutionally interpreted regions of a continuous,
path-dependent configuration space, rather than as discrete natural kinds
that a sufficiently powerful clustering algorithm would eventually recover
--- a reading the dissertation's own methodology, a coordinate space plus a
qualitative typology rather than a cluster analysis, is already consistent
with, and one this project's own generative and historical mechanisms are
jointly sufficient, though not yet calibrated to real historical data, to
produce and causally probe. Explaining why real governance systems
concentrate around particular recurring trajectories is deferred to future
historically calibrated research.
